\chapter{Type theory, the fixture}

\section{Homotopies}
\label{sec:homotopy}

\begin{defn}\label{defn:homotopy}
For functions $f,g:\prd{x:A}P(x)$, a \emph{homotopy} from $f$ to $g$ is a
function of the type $\prd{x:A}(f(x)=g(x))$.
\end{defn}

\begin{lem}\label{lem:homotopy-refl}
Each function has a homotopy to itself, in the sense of \cref{defn:homotopy}.
\end{lem}
\begin{proof}
Take the reflexivity path at each point.
\end{proof}

\begin{eg}
The identity function on $\nat$ has a homotopy to itself.
\end{eg}

\begin{rmk}
A remark is not an item.
\end{rmk}

\section*{Exercises}

\begin{ex}\label{ex:homotopy-sym}
Show that a homotopy from $f$ to $g$ gives one from $g$ to $f$, and use
\cref{lem:homotopy-refl,defn:homotopy}.
\end{ex}

% An example that the author did not write yet.
\begin{example}\label{eg:later}
\end{example}
